\usepackage{xcolor}
\usepackage{afterpage}
\usepackage{pifont,mdframed}
\usepackage[bottom]{footmisc}

\makeatletter
\gdef\this@inputfilename{input}
\gdef\this@outputfilename{output}
\makeatother

\lstset{
  basicstyle=\ttfamily,
  columns=fullflexible
}

\setcounter{@subtaskcounter}{-1}

%%%%%%%%%%%%%%%%%%%%%%%%%%%%%%%%%%%%%%%%%%%%%%%%%%%%%%%%%%%%%%%
%%%%%%%%%%%%%%%%%%%%%%%%%%%%%%%%%%%%%%%%%%%%%%%%%%%%%%%%%%%%%%%

%\textit{This problem is loosely based on real events.}
\textit{A feladat megtörtént eseményeken alapszik.}

%The organising committee of an unspecified programming contest (UPC) is planning to hold many rounds this year,
Egy meg nem nevezett programozási verseny szervezőbizottsága idén több forduló megrendezését tervezi,
%each round consisting of $N$ problems of varying difficulties. The problems in a round are numbered from $1$ to $N$.
melyek mindegyike $N$ darab, különböző nehézségű feladatból áll. Azt szeretnénk, hogy a feladatok nehézségi szintjei $1, 2, \ldots, N$ legyenek.

%As such, there have been a lot of problem proposals.
A versenyre rengeteg feladatjavaslat érkezett.
%Each problem proposal has a difficulty rating, which is either
Ezek mindegyikéhez tartozik egy nehézségi szint, mely lehet
\begin{itemize}
%\item an integer $i$ ($1 \le i \le N$) -- meaning that this problem is appropriate as problem $i$ in a round; or
\item egy $i$ egész szám ($1 \le i \le N$), ami azt jelenti, hogy a javaslat valamelyik forduló $i$ nehézségű feladataként használható; vagy
%\item a value in the form $i.5$ for some integer $i$ ($1 \le i < N$) -- meaning that this problem is appropriate as either problem $i$ or problem $i+1$ in a round.
\item egy $i.5$ alakú érték, ahol $i$ egész szám, $1 \le i < N$), mely akár $i$ nehézségű, akár $i+1$ nehézségű feladat is lehet.
\end{itemize}

\begin{figure}[H]
    \centering
    \includegraphics[width=0.6\linewidth]{problemsetting.jpg}
    %\caption{Creating a balanced round is more difficult than you'd think.}
    \caption{Egy megfelelő nehézségű verseny összeállítása nehezebb, mint gondolnád.}
\end{figure}

%Given the value of $N$ and the number of problem proposals for each difficulty rating,
%find the maximum number of rounds that can be created using these problems.
Az $N$ értékének és az egyes feladatok nehézségének ismeretében határozd meg, hogy legfeljebb hány fordulót lehet összeállítani ezen feladatokból.
%Each problem can only be used in at most one round.
Minden feladatot legfeljebb egy fordulóban használhatsz fel!

\begin{warning}
%Among the attachments of this task you may find a template file \texttt{problemsetting.*} with a sample incomplete implementation.
Az értékelő rendszerből letölthető csatolmányok közt találhatsz
\texttt{problemsetting.*} nevű fájlokat, melyek a bemeneti adatok beolvasását valósítják meg az egyes programnyelveken.
A megoldásodat ezekből a hiányos minta implementációkból kiindulva is elkészítheted.
\end{warning}



\InputFile

%Each test contains multiple test cases. The first line of input contains a single integer $T$, the number of test cases.
Minden teszt több tesztesetet tartalmaz. A bemenet első sora egyetlen $T$ egész számot tartalmaz, a tesztesetek számát.

%The first line of each test case contains the integer $N$, the number of problems in an UPC round.
Minden teszteset első sora egyetlen $N$ egész számot tartalmaz, az egy fordulóba választandó feladatok számát.

%The second line of each test case contains $N$ integers $A_1, A_2,\ldots, A_N$, where $A_i$ denotes the number of problem proposals with a difficulty rating of $i$.
A teszteset második sora $N$ darab,  $A_1, A_2,\ldots, A_N$ egész számot tartalmaz, ahol $A_i$ az $i$ nehézségű feladatok száma.

%The third line of each test case contains $N-1$ integers $B_1,B_2,\ldots,B_{N-1}$, where $B_i$ denotes the number of problem proposals with a difficulty rating of $i.5$.
A harmadik sorában $N-1$ darab $B_1,B_2,\ldots,B_{N-1}$ egész szám található, ahol $B_i$ az $i.5$ nehézségű feladatok száma.

\OutputFile

%For each test case, print a single integer, the maximum number of rounds that can be created from the available problems.
Minden tesztesethez egyetlen egész számot adj meg, az adott nehézségű feladatokból összeállítható fordulók maximális számát.


\Constraints
\begin{itemize}[nolistsep, itemsep=2mm]
    \item $1 \le T \le \constraints{MAXT}$.
	\item $2 \le N \le \constraints{MAXN}$.
	\item $0 \le A_i \le \constraints{MAXA}$ minden $i=1 \ldots N$-re.
	\item $0 \le B_i \le \constraints{MAXB}$ minden $i=1 \ldots N-1$-re.
    %\item The sum of $N$ across all test cases does not exceed $\constraints{MAXN}$.
    \item Az $N$-ek összege tesztenként nem haladja meg a $\constraints{MAXN}$-t.
\end{itemize}

\Scoring
A megoldásodat sok különböző tesztesetre lefuttatjuk. A tesztesetek részfeladatokba vannak csoportosítva. Egy-egy részfeladatot akkor tekintünk megoldottnak, ha volt legalább egy olyan beadásod, amely az adott részfeladat minden tesztesetére helyes megoldást adott.
A feladat összpontszámát a megoldott részfeladatokra kapott pontszámok összege adja.
%Your program will be tested against several test cases grouped in subtasks.
%In order to obtain the score of a subtask, your program needs to correctly solve all of its test cases.

\IIOTsubtask{0}{1}{Példák.}

\IIOTsubtask{5}{1}{$B_i=0$ minden $i=1 \ldots N-1$-re.}

\IIOTsubtask{10}{1}{$A_i \le 1$ minden $i=1 \ldots N$-re; $B_i \le 1$ minden $i=1 \ldots N-1$-re.}

\IIOTsubtask{30}{2}{$A_i \le 100$ minden $i=1 \ldots N$-re; $B_i \le 100$ minden $i=1 \ldots N-1$-re.}

\IIOTsubtask{20}{2}{$A_i=0$ minden $i=1 \ldots N$-re.}

\IIOTsubtask{35}{2}{Nincsenek további megkötések.}


\Examples
\begin{example}
\exmpfile{problemsetting.input0.txt}{problemsetting.output0.txt}%
\end{example}


\Explanation

%In the first test case, there are $7$ problem proposals, which have difficulty ratings of $1.5$, $2$, $3$, $3.5$, $5$, $5.5$ and $6$, respectively. 
Az \textbf{első tesztesetben} $7$ feladatjavaslat van, amelyeknek a nehézsége rendre $1.5$, $2$, $3$, $3.5$, $5$, $5.5$ és $6$.

%We can create a single round using these problems:
Ezekből a feladatokból egyetlen fordulót hozhatunk létre:
\begin{itemize}
%\item The problem with a difficulty rating of $1.5$ will be used as the first problem in the round.
\item A $1.5$-es nehézségi szintű javaslat a forduló első feladata lesz.

%\item The problem with a difficulty rating of $2$ will be used as the second problem in the round.
\item A $2$-es nehézségi szintű javaslat a forduló második feladata lesz.

%\item The problem with a difficulty rating of $3$ will be used as the third problem in the round.
\item A $3$-as nehézségi szintű javaslat a forduló harmadik feladata lesz.

%\item The problem with a difficulty rating of $3.5$ will be used as the fourth problem in the round.
\item A $3.5$-es nehézségi szintű javaslat a forduló negyedik feladata lesz.

%\item The problem with a difficulty rating of $5.5$ will be used as the fifth problem in the round.
\item A $5.5$-es nehézségi szintű javaslat a forduló ötödik feladata lesz.

%\item The problem with a difficulty rating of $6$ will be used as the sixth problem in the round.
\item A $6$-os nehézségi szintű javaslat a forduló hatodik feladata lesz.

\end{itemize}

%In the second test case, we can show that it is impossible to create a complete round from the given proposals.
Belátható, hogy nem hozható létre egyetlen teljes forduló sem a \textbf{második tesztesetben} található feladatokból.


%In the third test case, it is possible to create the following rounds from the given proposals:
A \textbf{harmadik tesztesetben} két fordulót állíthatunk össze, melyek ilyen nehézségű feladatokat tartalmaznak:
\begin{itemize}
\item $[1,1.5]$
\item $[1.5,2]$
\end{itemize}

%In the fourth test case, it is possible to create the following rounds from the given proposals:
A \textbf{negyedik tesztesetben} az alábbi fordulók létrehozása lehetséges:
\begin{itemize}
\item $[1.5,1.5,3.5,3.5,5.5,5.5]$
\item $[1.5,1.5,3.5,3.5,5.5,5.5]$
\item $[1.5,2.5,2.5,3.5,5.5,5.5]$
\end{itemize}
